% ============================================================
% Chapitre 7 : Tests Classiques z, t, chi2, F
% ============================================================
\chapter{Tests Classiques : $z$, $t$, $\chi^2$, $F$}
\label{ch:tests-classiques}

\begin{center}
\textit{<<~La th\'eorie sans la pratique est st\'erile ; la pratique sans la th\'eorie est aveugle.~>>}
\end{center}

% ------------------------------------------------------------
\section{Exemple motivant}
% ------------------------------------------------------------

Deux machines A et B produisent des pi\`eces. On veut comparer leur pr\'ecision (variance) et savoir si elles produisent en moyenne la m\^eme longueur. Pour la machine A : $n_1=20$, $\bar{x}_1=50{,}2$, $s_1=1{,}5$. Pour B : $n_2=25$, $\bar{x}_2=49{,}5$, $s_2=2{,}1$. Quels tests utiliser ?

% ------------------------------------------------------------
\section{Test $z$ (variance connue)}
\label{sec:test-z}
% ------------------------------------------------------------

\begin{theoreme}[Test $z$ pour la moyenne]
\index{test z@test $z$}
$X_1,\ldots,X_n\iid\mathcal{N}(\mu,\sigma^2)$, $\sigma^2$ connu. Pour $H_0:\mu=\mu_0$ :
\[
Z = \frac{\xbar-\mu_0}{\sigma/\sqrt{n}} \sim \mathcal{N}(0,1) \text{ sous } H_0.
\]
\begin{center}
\begin{tabular}{lll}
\toprule
$H_1$ & R\'egion critique & $p$-valeur \\
\midrule
$\mu\neq\mu_0$ & $|Z|>z_{\alpha/2}$ & $2\Phi(-|z_{\mathrm{obs}}|)$ \\
$\mu>\mu_0$ & $Z>z_\alpha$ & $1-\Phi(z_{\mathrm{obs}})$ \\
$\mu<\mu_0$ & $Z<-z_\alpha$ & $\Phi(z_{\mathrm{obs}})$ \\
\bottomrule
\end{tabular}
\end{center}
\end{theoreme}

\subsection{Test $z$ pour deux moyennes}

$X_1,\ldots,X_{n_1}\iid\mathcal{N}(\mu_1,\sigma_1^2)$ et $Y_1,\ldots,Y_{n_2}\iid\mathcal{N}(\mu_2,\sigma_2^2)$ ind\'ependants, variances connues :
\[
Z = \frac{\xbar-\bar{Y}-\delta_0}{\sqrt{\sigma_1^2/n_1+\sigma_2^2/n_2}}, \quad H_0:\mu_1-\mu_2=\delta_0.
\]

\subsection{Test $z$ pour une proportion}

$\phat=\xbar$ avec $X_i\iid\mathcal{B}(1,p)$, $n$ grand :
\[
Z = \frac{\phat-p_0}{\sqrt{p_0(1-p_0)/n}}, \quad H_0:p=p_0.
\]

\subsection{Test $z$ pour deux proportions}

$H_0:p_1=p_2$ :
\[
Z = \frac{\phat_1-\phat_2}{\sqrt{\phat_{\mathrm{pool}}(1-\phat_{\mathrm{pool}})(1/n_1+1/n_2)}}, \quad \phat_{\mathrm{pool}}=\frac{n_1\phat_1+n_2\phat_2}{n_1+n_2}.
\]

% ------------------------------------------------------------
\section{Test $t$ de Student}
\label{sec:test-t}
% ------------------------------------------------------------

\begin{theoreme}[Test $t$ pour un \'echantillon]
\index{test t@test $t$}
$X_1,\ldots,X_n\iid\mathcal{N}(\mu,\sigma^2)$, $\sigma^2$ inconnu. Pour $H_0:\mu=\mu_0$ :
\[
T = \frac{\xbar-\mu_0}{S/\sqrt{n}} \sim t(n-1) \text{ sous } H_0.
\]
R\'egions critiques analogues au test $z$ avec les quantiles de $t(n-1)$.
\end{theoreme}

\subsection{Test $t$ pour deux \'echantillons ind\'ependants}

\textbf{Variances \'egales :}
\[
T = \frac{\xbar-\bar{Y}}{S_p\sqrt{1/n_1+1/n_2}} \sim t(n_1+n_2-2), \quad S_p^2=\frac{(n_1-1)S_1^2+(n_2-1)S_2^2}{n_1+n_2-2}.
\]

\textbf{Variances in\'egales (Welch) :}
\[
T = \frac{\xbar-\bar{Y}}{\sqrt{S_1^2/n_1+S_2^2/n_2}} \sim t(\nu),
\]
avec $\nu$ calcul\'e par la formule de Welch--Satterthwaite :
\[
\nu = \frac{(S_1^2/n_1+S_2^2/n_2)^2}{\frac{(S_1^2/n_1)^2}{n_1-1}+\frac{(S_2^2/n_2)^2}{n_2-1}}.
\]

\subsection{Test $t$ appari\'e}

Pour des donn\'ees appari\'ees $(X_i,Y_i)$, on travaille sur les diff\'erences $D_i=X_i-Y_i$ :
\[
T = \frac{\bar{D}}{S_D/\sqrt{n}} \sim t(n-1), \quad H_0:\mu_D=0.
\]

% ------------------------------------------------------------
\section{Test du $\chi^2$}
\label{sec:test-chi2}
% ------------------------------------------------------------

\subsection{Test sur la variance}

\begin{theoreme}[Test $\chi^2$ pour la variance]
\index{test chi2@test $\chi^2$}
$X_1,\ldots,X_n\iid\mathcal{N}(\mu,\sigma^2)$. Pour $H_0:\sigma^2=\sigma_0^2$ :
\[
\chi^2 = \frac{(n-1)S^2}{\sigma_0^2} \sim \chi^2(n-1) \text{ sous } H_0.
\]
\end{theoreme}

\subsection{Test d'ad\'equation du $\chi^2$}

\begin{theoreme}[Test d'ad\'equation de Pearson]
\index{test d'adequation@test d'ad\'equation}
\label{thm:chi2-adequation}
Soit $k$ cat\'egories avec effectifs observ\'es $O_i$ et effectifs th\'eoriques $E_i=np_i$ :
\[
\chi^2 = \sum_{i=1}^k \frac{(O_i-E_i)^2}{E_i} \xrightarrow{\mathcal{L}} \chi^2(k-1-r),
\]
o\`u $r$ est le nombre de param\`etres estim\'es. On rejette $H_0$ pour $\chi^2>\chi^2_{\alpha,k-1-r}$.
\end{theoreme}

\subsection{Test d'ind\'ependance du $\chi^2$}

\begin{theoreme}[Test d'ind\'ependance]
\index{test d'independance@test d'ind\'ependance}
Pour un tableau de contingence $r\times c$ avec effectifs $O_{ij}$ :
\[
\chi^2 = \sum_{i=1}^r\sum_{j=1}^c \frac{(O_{ij}-E_{ij})^2}{E_{ij}}, \quad E_{ij}=\frac{O_{i\cdot}\cdot O_{\cdot j}}{n},
\]
suit approximativement $\chi^2((r-1)(c-1))$ sous $H_0$ : ind\'ependance.
\end{theoreme}

% ------------------------------------------------------------
\section{Test $F$ de Fisher}
\label{sec:test-f}
% ------------------------------------------------------------

\begin{theoreme}[Test $F$ pour le rapport de variances]
\index{test F@test $F$}
$X_1,\ldots,X_{n_1}\iid\mathcal{N}(\mu_1,\sigma_1^2)$ et $Y_1,\ldots,Y_{n_2}\iid\mathcal{N}(\mu_2,\sigma_2^2)$ ind\'ependants. Pour $H_0:\sigma_1^2=\sigma_2^2$ :
\[
F = \frac{S_1^2}{S_2^2} \sim F(n_1-1,n_2-1) \text{ sous } H_0.
\]
On rejette $H_0$ si $F>F_{\alpha/2,n_1-1,n_2-1}$ ou $F<F_{1-\alpha/2,n_1-1,n_2-1}$.
\end{theoreme}

% ------------------------------------------------------------
\section{R\'esum\'e : choix du test}
% ------------------------------------------------------------

\begin{center}
\begin{tikzpicture}[
  box/.style={rectangle, draw, rounded corners, minimum width=3cm, minimum height=1cm, text centered, font=\small},
  arrow/.style={-{Stealth}, thick}
]
\node[box, fill=blue!10] (start) at (0,0) {Param\`etre ?};
\node[box, fill=green!10] (mu) at (-4,-2) {Moyenne $\mu$};
\node[box, fill=green!10] (sig) at (0,-2) {Variance $\sigma^2$};
\node[box, fill=green!10] (prop) at (4,-2) {Proportion $p$};
\node[box, fill=orange!10] (z) at (-6,-4) {$\sigma$ connu $\to$ $z$};
\node[box, fill=orange!10] (t) at (-2,-4) {$\sigma$ inconnu $\to$ $t$};
\node[box, fill=orange!10] (chi) at (0,-4) {$\chi^2$};
\node[box, fill=orange!10] (zprop) at (4,-4) {$z$ (grands $n$)};

\draw[arrow] (start) -- (mu);
\draw[arrow] (start) -- (sig);
\draw[arrow] (start) -- (prop);
\draw[arrow] (mu) -- (z);
\draw[arrow] (mu) -- (t);
\draw[arrow] (sig) -- (chi);
\draw[arrow] (prop) -- (zprop);
\end{tikzpicture}
\end{center}

% ------------------------------------------------------------
\section{Impl\'ementation Python}
% ------------------------------------------------------------

\begin{python}
\begin{minted}{python}
import numpy as np
from scipy import stats

# --- Test t pour un echantillon ---
data = np.array([50.2, 49.8, 51.1, 48.9, 50.5, 49.3, 50.8, 49.1,
                 50.0, 49.7, 50.3, 49.5, 50.6, 49.2, 50.4])
t_stat, p_val = stats.ttest_1samp(data, 50)
print(f"Test t (H0: mu=50): t={t_stat:.4f}, p={p_val:.4f}")

# --- Test t pour deux echantillons ---
A = np.array([50.2, 49.8, 51.1, 48.9, 50.5, 49.3, 50.8, 49.1,
              50.0, 49.7, 50.3, 49.5, 50.6, 49.2, 50.4,
              50.1, 49.6, 50.7, 49.4, 50.2])
B = np.array([49.5, 48.2, 50.1, 47.8, 49.0, 48.5, 50.3, 47.5,
              49.2, 48.8, 49.7, 48.0, 49.4, 48.3, 49.1,
              48.7, 49.6, 47.9, 49.3, 48.1, 49.8, 48.4,
              49.0, 48.6, 49.5])

# Variances egales
t_eq, p_eq = stats.ttest_ind(A, B, equal_var=True)
print(f"Test t (egal var): t={t_eq:.4f}, p={p_eq:.4f}")

# Welch (variances inegales)
t_w, p_w = stats.ttest_ind(A, B, equal_var=False)
print(f"Test Welch:        t={t_w:.4f}, p={p_w:.4f}")

# --- Test F pour egalite des variances ---
f_stat = np.var(A, ddof=1) / np.var(B, ddof=1)
df1, df2 = len(A)-1, len(B)-1
p_f = 2 * min(stats.f.cdf(f_stat, df1, df2), 1-stats.f.cdf(f_stat, df1, df2))
print(f"Test F: F={f_stat:.4f}, p={p_f:.4f}")

# --- Test chi2 d'adequation ---
# De de 60 lancers
obs = np.array([8, 12, 10, 11, 9, 10])  # observes
exp = np.array([10, 10, 10, 10, 10, 10])  # attendus (de equilibre)
chi2, p_chi = stats.chisquare(obs, exp)
print(f"chi2 adequation: chi2={chi2:.4f}, p={p_chi:.4f}")

# --- Test chi2 d'independance ---
# Tableau de contingence
tableau = np.array([[30, 20, 10],
                    [25, 30, 15],
                    [15, 20, 35]])
chi2_ind, p_ind, dof, expected = stats.chi2_contingency(tableau)
print(f"chi2 independance: chi2={chi2_ind:.4f}, p={p_ind:.4f}, ddl={dof}")

# --- Test de proportion ---
from statsmodels.stats.proportion import proportions_ztest
count = np.array([520])
nobs = np.array([1000])
z_prop, p_prop = proportions_ztest(count, nobs, value=0.5)
print(f"Test proportion (H0: p=0.5): z={z_prop:.4f}, p={p_prop:.4f}")
\end{minted}
\end{python}

% ------------------------------------------------------------
\section{Exercices}
% ------------------------------------------------------------

\begin{exercice}[$\star$ -- Tests sur la moyenne]
Les temps de r\'eaction (en ms) de 12 sujets apr\`es un caf\'e sont : 210, 195, 220, 205, 198, 215, 208, 200, 225, 192, 218, 202. Le temps moyen normal est 220~ms. Le caf\'e r\'eduit-il significativement le temps de r\'eaction ? (Test $t$ unilat\'eral, $\alpha=0{,}05$.)
\end{exercice}

\begin{exercice}[$\star$ -- Test $\chi^2$]
On lance un d\'e 120 fois et on observe : face 1 : 25, face 2 : 17, face 3 : 22, face 4 : 18, face 5 : 19, face 6 : 19. Le d\'e est-il \'equilibr\'e ? ($\alpha=0{,}05$.)
\end{exercice}

\begin{exercice}[$\star\star$ -- Deux \'echantillons]
Comparer les moyennes de deux traitements : Traitement A ($n_1=15$, $\bar{x}_1=72{,}3$, $s_1=8{,}5$) et Traitement B ($n_2=18$, $\bar{x}_2=68{,}1$, $s_2=9{,}2$). D'abord tester l'\'egalit\'e des variances (test $F$), puis choisir le test $t$ adapt\'e.
\end{exercice}

\begin{exercice}[$\star\star$ -- Test appari\'e]
Dix patients ont leur pression art\'erielle mesur\'ee avant et apr\`es un traitement. Avant : 140, 135, 150, 145, 138, 142, 155, 148, 137, 143. Apr\`es : 132, 130, 142, 138, 135, 136, 148, 140, 133, 138. Le traitement est-il efficace ?
\end{exercice}

\begin{exercice}[$\star\star\star$ -- Projet : robustesse des tests]
Par simulation, \'etudier la robustesse du test $t$ \`a la non-normalit\'e. G\'en\'erer des \'echantillons de lois exponentielle, log-normale et uniforme. Pour $n=5,10,30,100$, comparer le taux d'erreur de type I r\'eel au niveau nominal 5\%.
\end{exercice}

% ------------------------------------------------------------
\begin{formules}
\begin{align*}
z &= \frac{\xbar-\mu_0}{\sigma/\sqrt{n}} & t &= \frac{\xbar-\mu_0}{S/\sqrt{n}} \\
\chi^2 &= \frac{(n-1)S^2}{\sigma_0^2} & F &= \frac{S_1^2}{S_2^2} \\
\chi^2_{\text{ad\'eq}} &= \sum\frac{(O_i-E_i)^2}{E_i} & \chi^2_{\text{ind}} &= \sum\frac{(O_{ij}-E_{ij})^2}{E_{ij}} \\
S_p^2 &= \frac{(n_1-1)S_1^2+(n_2-1)S_2^2}{n_1+n_2-2} & z_{\text{prop}} &= \frac{\phat-p_0}{\sqrt{p_0(1-p_0)/n}}
\end{align*}
\end{formules}
